\chapter{2020年北京卷}

\section{单选题}

\begin{problem}
已知集合 \(A=\left\{ -1,0,1,2 \right\}\)，\(B=\left\{ x\mid 0<x<3 \right\}\)，则 \(A\cap B=\)（\quad）.

\choices
    {\(\left\{ -1,0,1 \right\}\)}
    {\(\left\{ 0,1 \right\}\)}
    {\(\left\{ -1,1,2 \right\}\)}
    {\(\left\{ 1,2 \right\}\)}

\end{problem}

\begin{answer}
D
\end{answer}

\begin{solution}
由交集的定义可得
\(A\cap B=\left\{ -1,0,1,2 \right\}\cap \left\{ x\mid 0<x<3 \right\}=\left\{ 1,2 \right\}\).
故选 D.
\end{solution}

\begin{problem}
在复平面内，复数 \(z\) 对应的点的坐标是 \((1,2)\)，则 \(\i z=\)（\quad）.

\choices
    {\(1+2\i\)}
    {\(-2+\i\)}
    {\(1-2\i\)}
    {\(-2-\i\)}

\end{problem}

\begin{answer}
B
\end{answer}

\begin{solution}
由题意得 \(z=1+2\i\)，所以
\(\i z=\i(1+2\i)=-2+\i\).
故选 B.
\end{solution}

\begin{problem}
在 \(\left(\sqrt{x}-2\right)^5\) 的展开式中，\(x^2\) 的系数为（\quad）.

\choices
    {\(-5\)}
    {\(5\)}
    {\(-10\)}
    {\(10\)}

\end{problem}

\begin{answer}
C
\end{answer}

\begin{solution}
二项展开式的通项为
\(T_{r+1}=\mathrm{C}_5^r(\sqrt{x})^{5-r}(-2)^r\)（\(r=0\)，\(1\)，\(2\)，\(3\)，\(4\)，\(5\)）.
令 \(\dfrac{5-r}{2}=2\)，得 \(r=1\). 因此 \(x^2\) 的系数为
\(\mathrm{C}_5^1(-2)=-10\).
故选 C.
\end{solution}

\begin{problem}
某三棱柱的底面为正三角形，其三视图如图所示，该三棱柱的表面积为（\quad）.

\begin{center}
\bitmapfigure[max width=0.42\linewidth]{img/2020/beijing/q4_fig1.png}
\end{center}

\choices
    {\(6+3\sqrt3\)}
    {\(6+2\sqrt3\)}
    {\(12+3\sqrt3\)}
    {\(12+2\sqrt3\)}

\end{problem}

\begin{answer}
D
\end{answer}

\begin{solution}
由三视图可知，三棱柱的上下底面为边长为 \(2\) 的等边三角形，侧面为三个边长为 \(2\) 的正方形，所以其表面积为
\(S=3\times(2\times2)+2\times\left(\frac12\times2\times2\times\sin60^\circ\right)=12+2\sqrt3\).
故选 D.
\end{solution}

\begin{problem}
已知半径为 \(1\) 的圆经过点 \((3,4)\)，则其圆心到原点的距离的最小值为（\quad）.

\choices
    {\(4\)}
    {\(5\)}
    {\(6\)}
    {\(7\)}

\end{problem}

\begin{answer}
A
\end{answer}

\begin{solution}
设圆心为 \(C\)，则 \(C\) 的轨迹是以 \(M(3,4)\) 为圆心，\(1\) 为半径的圆. 于是
\(|OC|+1\ge |OM|=\sqrt{3^2+4^2}=5\)，
所以
\(|OC|\ge 4\).
当且仅当 C 在线段 \(OM\) 上时取等号，故最小值为 \(4\). 故选 A.
\end{solution}

\begin{problem}
已知函数 \(f(x)=2^x-x-1\)，则不等式 \(f(x)>0\) 的解集是（\quad）.

\choices
    {\(( -1,1)\)}
    {\((-\infty,-1)\cup(1,+\infty)\)}
    {\((0,1)\)}
    {\((-\infty,0)\cup(1,+\infty)\)}

\end{problem}

\begin{answer}
D
\end{answer}

\begin{solution}
因为
\(f(x)=2^x-x-1\)，
所以 \(f(x)>0\) 等价于 \(2^x>x+1\). 在同一直角坐标系中作出 \(y=2^x\) 和 \(y=x+1\) 的图象. 两函数图象的交点坐标为 \((0,1)\)，\((1,2)\)，因此不等式 \(2^x>x+1\) 的解为 \(x<0\) 或 \(x>1\). 所以不等式 \(f(x)>0\) 的解集为 \((-\infty,0)\cup(1,+\infty)\). 故选 D.
\end{solution}

\begin{problem}
设抛物线的顶点为 \(O\)，焦点为 \(F\)，准线为 \(l\). \(P\) 是抛物线上异于 \(O\) 的一点，过 \(P\) 作 \(PQ\perp l\) 于 \(Q\)，则线段 \(FQ\) 的垂直平分线（\quad）.

\choices
    {经过点 \(O\)}
    {经过点 \(P\)}
    {平行于直线 \(OP\)}
    {垂直于直线 \(OP\)}

\end{problem}

\begin{answer}
B
\end{answer}

\begin{solution}
依据题意不妨作出焦点在 \(x\) 轴上的开口向右的抛物线. 根据垂直平分线的定义和抛物线的定义可知，线段 \(FQ\) 的垂直平分线经过点 \(P\). 故选 B.
\end{solution}

\begin{problem}
在等差数列 \(\{a_n\}\) 中，\(a_1=-9\)，\(a_3=-1\). 记 \(T_n=a_1a_2\cdots a_n\)（\(n=1\)，\(2\)，\(\dots\)），则数列 \(\{T_n\}\)（\quad）.

\choices
    {有最大项，有最小项}
    {有最大项，无最小项}
    {无最大项，有最小项}
    {无最大项，无最小项}

\end{problem}

\begin{answer}
B
\end{answer}

\begin{solution}
由题意可知等差数列的公差为
\(d=\frac{a_3-a_1}{3-1}=\frac{-1+9}{2}=2\).
所以 \(a_n=a_1+(n-1)d=-9+2(n-1)=2n-11\).
注意到 \(a_1<a_2<a_3<a_4<a_5<0<a_6=1<a_7<\cdots\).
且由 \(T_5<0\) 可知 \(T_i<0\) （\(i\ge6\)），由
\(\frac{T_i}{T_{i-1}}=a_i>1\) （\(i\ge7\)）可知数列 \(\{T_n\}\) 不存在最小项.
由于 \(a_1=-9\)，\(a_2=-7\)，\(a_3=-5\)，\(a_4=-3\)，\(a_5=-1\)，\(a_6=1\)，故数列 \(\{T_n\}\) 中的正项只有有限项：\(T_2=63\)，\(T_4=945\).
故数列 \(\{T_n\}\) 中存在最大项，且最大项为 \(T_4\). 故选 B.
\end{solution}

\begin{problem}
已知 \(\alpha\)，\(\beta\in\R\)，则“存在 \(k\in\Z\) 使得 \(\alpha=k\pi+(-1)^k\beta\)”是“\(\sin\alpha=\sin\beta\)”的（\quad）.

\choices
    {充分而不必要条件}
    {必要而不充分条件}
    {充分必要条件}
    {既不充分也不必要条件}

\end{problem}

\begin{answer}
C
\end{answer}

\begin{solution}
当存在 \(k\in\Z\) 使得 \(\alpha=k\pi+(-1)^k\beta\) 时，若 \(k\) 为偶数，则
\(\sin\alpha=\sin(k\pi+\beta)=\sin\beta\).
若 \(k\) 为奇数，则
\(\sin\alpha=\sin(k\pi-\beta)=\sin\left((k-1)\pi+\pi-\beta\right)=\sin(\pi-\beta)=\sin\beta\).
反过来，当 \(\sin\alpha=\sin\beta\) 时，\(\alpha=\beta+2m\pi\) 或 \(\alpha+\beta=\pi+2m\pi\)，\(m\in\Z\)，即
\(\alpha=k\pi+(-1)^k\beta\).
所以题述关系是充要条件. 故选 C.
\end{solution}

\begin{problem}
2020 年 3 月 14 日是全球首个国际圆周率日（\(\pi\) Day）. 历史上，求圆周率 \(\pi\) 的方法有多种，与中国传统数学中的“割圆术”相似. 数学家阿尔·卡西的方法是：当正整数 \(n\) 充分大时，计算单位圆的内接正 \(6n\) 边形的周长和外切正 \(6n\) 边形（各边均与圆相切的正 \(6n\) 边形）的周长，将它们的算术平均数作为 \(2\pi\) 的近似值. 按照阿尔·卡西的方法，\(\pi\) 的近似值的表达式是（\quad）.

\choices
    {\(3n\left(\sin\frac{30^\circ}{n}+\tan\frac{30^\circ}{n}\right)\)}
    {\(6n\left(\sin\frac{30^\circ}{n}+\tan\frac{30^\circ}{n}\right)\)}
    {\(3n\left(\sin\frac{60^\circ}{n}+\tan\frac{60^\circ}{n}\right)\)}
    {\(6n\left(\sin\frac{60^\circ}{n}+\tan\frac{60^\circ}{n}\right)\)}

\end{problem}

\begin{answer}
A
\end{answer}

\begin{solution}
单位圆内接正 \(6n\) 边形的每条边所对应的圆周角为
\(\frac{360^\circ}{6n}=\frac{60^\circ}{n}\)，
每条边长为
\(2\sin\frac{30^\circ}{n}\).
所以，单位圆内接正 \(6n\) 边形的周长为
\(12n\sin\frac{30^\circ}{n}\).
单位圆外切正 \(6n\) 边形的每条边长为
\(2\tan\frac{30^\circ}{n}\)，
其周长为
\(12n\tan\frac{30^\circ}{n}\).
于是 \(2\pi=\frac{12n\sin\frac{30^\circ}{n}+12n\tan\frac{30^\circ}{n}}{2}=6n\left(\sin\frac{30^\circ}{n}+\tan\frac{30^\circ}{n}\right)\)，所以 \(\pi=3n\left(\sin\frac{30^\circ}{n}+\tan\frac{30^\circ}{n}\right)\).
故选 A.
\end{solution}
\section{填空题}

\begin{problem}
函数 \(\dfrac{1}{x+1}+\ln x\) 的定义域是\fillinblank{}.

\end{problem}

\begin{answer}
\((0,+\infty)\)
\end{answer}

\begin{solution}
由题意得 \(x>0\)，\(x+1\ne0\)，
故定义域为 \((0,+\infty)\).
\end{solution}

\begin{problem}
已知双曲线 \(C:\dfrac{x^2}{6}-\dfrac{y^2}{3}=1\)，则 \(C\) 的右焦点的坐标为\fillinblank{}；\(C\) 的焦点到其渐近线的距离是\fillinblank{}.

\end{problem}

\begin{answer}
\((3,0)\)，\(\ \sqrt3\)
\end{answer}

\begin{solution}
在双曲线 \(C\) 中，\(a^2=6\)，\(b^2=3\)，所以
\(c^2=a^2+b^2=9\)，\(c=3\).
因此右焦点坐标为 \((3,0)\). 双曲线的渐近线方程为
\(y=\pm\frac{b}{a}x=\pm\frac{1}{\sqrt2}x\)，
即 \(x\pm \sqrt{2}y=0\)，所以焦点到渐近线的距离为 \(\sqrt3\). 故答案为 \((3,0)\) 和 \(\sqrt3\).
\end{solution}

\begin{problem}
已知正方形 \(ABCD\) 的边长为 \(2\)，点 \(P\) 满足 \(\overrightarrow{AP}=\dfrac12\left(\overrightarrow{AB}+\overrightarrow{AC}\right)\)，则 \(|PD|=\fillinblank{}\)；\(\overrightarrow{PB}\cdot\overrightarrow{PD}=\fillinblank{}\).

\end{problem}

\begin{answer}
\(\sqrt5\)，\(\ -1\)
\end{answer}

\begin{solution}
以点 \(A\) 为坐标原点，\(AB\)，\(AD\) 所在直线分别为 \(x\)，\(y\) 轴建立如下图所示的平面直角坐标系：

\solutionfigure{\bitmapfigure[max width=0.42\linewidth]{img/2020/beijing/q13_fig1.png}}

则
\(A(0,0)\)，\(B(2,0)\)，\(C(2,2)\)，\(D(0,2)\).
由 \(\overrightarrow{AP}=\frac12\left(\overrightarrow{AB}+\overrightarrow{AC}\right)=\frac12\left((2,0)+(2,2)\right)=(2,1)\)，得 \(P(2,1)\). 于是 \(\overrightarrow{PD}=(-2,1)\)，\(\overrightarrow{PB}=(0,-1)\).
所以 \(|PD|=\sqrt{(-2)^2+1^2}=\sqrt5\)，\(\overrightarrow{PB}\cdot\overrightarrow{PD}=0\cdot(-2)+(-1)\cdot1=-1\).
故答案为 \(\sqrt5\) 和 \(-1\).
\end{solution}

\begin{problem}
若函数 \(f(x)=\sin(x+\varphi)+\cos x\) 的最大值为 \(2\)，则常数 \(\varphi\) 的一个取值为\fillinblank{}.

\end{problem}

\begin{answer}
\(\dfrac{\pi}{2}\)
\end{answer}

\begin{solution}
因为
\(f(x)=\sin(x+\varphi)+\cos x=\cos\varphi\sin x+(\sin\varphi+1)\cos x\).
由辅助角公式可知，\(f(x)\) 的最大值为
\(\sqrt{\cos^2\varphi+(\sin\varphi+1)^2}\).
题设给出最大值为 \(2\)，于是
\(\cos^2\varphi+(\sin\varphi+1)^2=4\).
化简得
\(2\sin\varphi+2=4\)，
所以 \(\sin\varphi=1\)，可取 \(\varphi=\dfrac{\pi}{2}\).
\end{solution}

\begin{problem}
为满足人民对美好生活的向往，环保部门要求相关企业加强污水治理，排放未达标的企业要限期整改，设企业的污水排放量 \(W\) 与时间 \(t\) 的关系为 \(W=f(t)\)，用 \(-\dfrac{f(b)-f(a)}{b-a}\) 的大小评价在 \([a,b]\) 这段时间内企业污水治理能力的强弱，已知整改期内，甲，乙两企业的污水排放量与时间的关系如下图所示.

\begin{center}
\bitmapfigure[max width=0.80\linewidth]{img/2020/beijing/q15_fig1.png}
\end{center}

给出下列四个结论：

\begin{circlelist}
\item 在 \([t_1,t_2]\) 这段时间内，甲企业的污水治理能力比乙企业强；

\item 在 \(t_2\) 时刻，甲企业的污水治理能力比乙企业强；

\item 在 \(t_3\) 时刻，甲，乙两企业的污水排放都已达标；

\item 甲企业在 \([0,t_1]\)，\([t_1,t_2]\)，\([t_2,t_3]\) 这三段时间中，在 \([0,t_1]\) 的污水治理能力最强.
\end{circlelist}
其中所有正确结论的序号是\fillinblank{}.

\end{problem}

\begin{answer}
\circled{1}\circled{2}\circled{3}
\end{answer}

\begin{solution}
由题意，依次比较各区间端点连线斜率的负数，以及 \(t_2\) 时刻两条切线斜率的负数.
\(-\frac{f(b)-f(a)}{b-a}\) 表示区间端点连线斜率的负数.
在 \([t_1,t_2]\) 这段时间内，甲的斜率比乙的小，所以甲的斜率的相反数比乙的大，因此甲企业的污水治理能力比乙企业强，\circled{1}正确.

甲企业在 \([0,t_1]\)，\([t_1,t_2]\)，\([t_2,t_3]\) 这三段时间中，甲企业在 \([t_1,t_2]\) 这段时间内斜率最小，其相反数最大，即在 \([t_1,t_2]\) 的污水治理能力最强，\circled{4}错误.

在 \(t_2\) 时刻，甲切线的斜率比乙的小，所以甲切线的斜率的相反数比乙的大，甲企业的污水治理能力比乙企业强，\circled{2}正确.

在 \(t_3\) 时刻，甲，乙两企业的污水排放量都在污水达标排放量以下，所以都已达标，\circled{3}正确.

故答案为：\circled{1}\circled{2}\circled{3}.
\end{solution}
\section{解答题}

\begin{problem}
如图，在正方体 \(ABCD-A_1B_1C_1D_1\) 中，\(E\) 为 \(BB_1\) 的中点.



\begin{enumerate}
    \item 求证：\(BC_1\parallel\) 平面 \(AD_1E\)；
    \item 求直线 \(AA_1\) 与平面 \(AD_1E\) 所成角的正弦值.
\end{enumerate}

\begin{center}
\bitmapfigure[max width=0.60\linewidth]{img/2020/beijing/q16_fig1.png}
\end{center}

\end{problem}

\begin{solution}
【分析】
（1）先证明四边形 \(ABC_1D_1\) 为平行四边形，从而得到 \(BC_1\parallel AD_1\)；再由 \(BC_1\not\subset\) 平面 \(AD_1E\)，\(AD_1\subset\) 平面 \(AD_1E\)，说明 \(BC_1\parallel\) 平面 \(AD_1E\)；

（2）以点 \(A\) 为坐标原点，\(AD\)，\(AB\)，\(AA_1\) 所在直线分别为 \(x\)，\(y\)，\(z\) 轴建立空间直角坐标系 \(A-xyz\)，利用空间向量法可计算出直线 \(AA_1\) 与平面 \(AD_1E\) 所成角的正弦值.

【详解】

（1）如下图所示：

\solutionfigure{\bitmapfigure[max width=0.60\linewidth]{img/2020/beijing/q16_solution_fig1.png}}

在正方体 \(ABCD-A_1B_1C_1D_1\) 中，\(AB\parallel A_1B_1\) 且 \(AB=A_1B_1\)，\(A_1B_1\parallel C_1D_1\) 且 \(A_1B_1=C_1D_1\)，所以 \(AB\parallel C_1D_1\) 且 \(AB=C_1D_1\)，所以四边形 \(ABC_1D_1\) 为平行四边形，则 \(BC_1\parallel AD_1\).
又 \(BC_1\not\subset\) 平面 \(AD_1E\)，\(AD_1\subset\) 平面 \(AD_1E\)，所以 \(BC_1\parallel\) 平面 \(AD_1E\).

（2）以点 \(A\) 为坐标原点，\(AD\)，\(AB\)，\(AA_1\) 所在直线分别为 \(x\)，\(y\)，\(z\) 轴建立空间直角坐标系 \(A-xyz\)，如下图所示：

\solutionfigure{\bitmapfigure[max width=0.60\linewidth]{img/2020/beijing/q16_solution_fig2.png}}

设正方体 \(ABCD-A_1B_1C_1D_1\) 的棱长为 \(2\)，则
\(A(0,0,0)\)，\(A_1(0,0,2)\)，\(D_1(2,0,2)\)，\(E(0,2,1)\).
于是
\(\overrightarrow{AD_1}=(2,0,2)\)，\(\overrightarrow{AE}=(0,2,1)\).
设平面 \(AD_1E\) 的法向量为 \(\bs n=(x,y,z)\)，则
\(\bs n\cdot\overrightarrow{AD_1}=0\)，\(\bs n\cdot\overrightarrow{AE}=0\)，
即
\(2x+2z=0\)，\(2y+z=0\).
令 \(z=-2\)，得 \(x=2\)，\(y=1\)，所以 \(\bs n=(2,1,-2)\).
因此，直线 \(AA_1\) 与平面 \(AD_1E\) 所成角的正弦值为 \(\sin\theta=\frac{|\overrightarrow{AA_1}\cdot\bs n|}{|\overrightarrow{AA_1}||\bs n|}=\frac{4}{2\cdot3}=\frac{2}{3}\).
\end{solution}

\begin{problem}
在 \(\triangle ABC\) 中，\(a+b=11\)，再从条件（1），条件（2）这两个条件中选择一个作为已知，求：

\begin{enumerate}
    \item \(a\) 的值；
    \item \(\sin C\) 和 \(\triangle ABC\) 的面积.
\end{enumerate}

条件（1）：\(c=7\)，\(\cos A=-\dfrac17\)；

条件（2）：\(\cos A=\dfrac18\)，\(\cos B=\dfrac{9}{16}\).

\end{problem}

\begin{solution}
选择条件（1）.

（1）由余弦定理得
\(a^2=b^2+c^2-2bc\cos A\).
又 \(a+b=11\)，\(c=7\)，\(\cos A=-\dfrac17\)，所以
\(a^2=(11-a)^2+7^2-2(11-a)\cdot7\cdot\left(-\frac17\right)\)，
解得 \(a=8\).

（2）由
\(\cos A=-\frac17\)，\(A\in(0,\pi)\)，
得
\(\sin A=\sqrt{1-\cos^2A}=\frac{4\sqrt3}{7}\).
由正弦定理得
\(\frac{a}{\sin A}=\frac{c}{\sin C}\)，
即
\(\frac{8}{4\sqrt3/7}=\frac{7}{\sin C}\)，
解得
\(\sin C=\frac{\sqrt3}{2}\).
于是
\(S=\frac12ab\sin C=\frac12(11-8)\cdot8\cdot\frac{\sqrt3}{2}=6\sqrt3\).

选择条件（2）.

（1）由
\(\cos A=\frac18\)，\(\cos B=\frac{9}{16}\)，
得
\(\sin A=\frac{3\sqrt7}{8}\)，\(\sin B=\frac{5\sqrt7}{16}\).
由正弦定理得
\(\frac{a}{\sin A}=\frac{b}{\sin B}\)，
即
\(\frac{a}{3\sqrt7/8}=\frac{11-a}{5\sqrt7/16}\)，
解得 \(a=6\).

（2）由两角和正弦公式得
\(\sin C=\sin(A+B)=\sin A\cos B+\sin B\cos A\)
\(=\frac{3\sqrt7}{8}\cdot\frac{9}{16}+\frac{5\sqrt7}{16}\cdot\frac18 =\frac{\sqrt7}{4}\).
于是
\(S=\frac12ab\sin C=\frac12\cdot6\cdot5\cdot\frac{\sqrt7}{4}=\frac{15\sqrt7}{4}\).
\end{solution}

\begin{problem}
某校为举办甲，乙两项不同活动，分别设计了相应的活动方案：方案一，方案二. 为了解该校学生对活动方案是否支持，对学生进行简单随机抽样，获得数据如下表：

\begin{center}
\begin{tabular}{c|cc|cc}
\hline
& \multicolumn{2}{c|}{男生} & \multicolumn{2}{c}{女生} \\
\hline
& 支持 & 不支持 & 支持 & 不支持 \\
\hline
方案一 & \(200\text{ 人}\) & \(400\text{ 人}\) & \(300\text{ 人}\) & \(100\text{ 人}\) \\
方案二 & \(350\text{ 人}\) & \(250\text{ 人}\) & \(150\text{ 人}\) & \(250\text{ 人}\) \\
\hline
\end{tabular}
\end{center}

假设所有学生对活动方案是否支持相互独立.

\begin{enumerate}
    \item 分别估计该校男生支持方案一的概率，该校女生支持方案一的概率；
    \item 从该校全体男生中随机抽取 \(2\text{ 人}\)，全体女生中随机抽取 \(1\text{ 人}\)，估计这 \(3\text{ 人}\)中恰有 \(2\text{ 人}\)支持方案一的概率；
    \item 将该校学生支持方案一的概率估计值记为 \(p_0\)，假设该校年级有 \(500\text{ 名}\) 男生和 \(300\text{ 名}\) 女生，除一年级外其他年级学生支持方案二的概率估计值记为 \(p_1\)，试比较 \(p_0\) 与 \(p_1\) 的大小. （结论不要求证明）
\end{enumerate}

\end{problem}

\begin{solution}
（1）该校男生支持方案一的概率为
\(\frac{200}{200+400}=\frac13\)，
该校女生支持方案一的概率为
\(\frac{300}{300+100}=\frac34\).

（2）\(3\text{ 人}\) 中恰有 \(2\text{ 人}\) 支持方案一分两种情况：\(1^{\odot}\) 仅有两个男生支持方案一，\(2^{\odot}\) 仅有一个男生支持方案一，一个女生支持方案一，所以 \(P=\left(\frac13\right)^2\left(1-\frac34\right)+\mathrm C_2^1\cdot\frac13\left(1-\frac13\right)\cdot\frac34=\frac{13}{36}\).

（3）由题意可得
\(p_0=\frac{500\times\frac13+300\times\frac34}{500+300}=\frac{47}{96}\)，
\(p_1=\frac{500\times\frac{350}{350+250}+300\times\frac{150}{150+250}}{500+300}=\frac{97}{192}\).
所以 \(p_0<p_1\).
\end{solution}

\begin{problem}
已知函数 \(f(x)=12-x^2\).

\begin{enumerate}
    \item 求曲线 \(y=f(x)\) 的斜率等于 \(-2\) 的切线方程；
    \item 设曲线 \(y=f(x)\) 在点 \((t,f(t))\) 处的切线与坐标轴围成的三角形的面积为 \(S(t)\)，求 \(S(t)\) 的最小值.
\end{enumerate}

\end{problem}

\begin{solution}
（1）因为
\(f'(x)=-2x\)，
所以设切点为 \((x_0,12-x_0^2)\)，则
\(-2x_0=-2\)，
即 \(x_0=1\)，所以切点为 \((1,11)\). 由点斜式可得切线方程
\(y-11=-2(x-1)\)，
即
\(2x+y-13=0\).

（2）曲线 \(y=f(x)\) 在点 \((t,f(t))\) 处的切线方程为
\(y-\left(12-t^2\right)=-2t(x-t)\).
令 \(x=0\)，得 \(y=t^2+12\)；令 \(y=0\)，得
\(x=\frac{t^2+12}{2t}\).
所以
\(S(t)=\frac12\cdot (t^2+12)\cdot \frac{t^2+12}{2|t|} =\frac{(t^2+12)^2}{4|t|}\).
不妨设 \(t>0\)（\(t<0\) 时结果相同），则
\(S(t)=\frac{t^4+24t^2+144}{4t} =\frac14\left(t^3+24t+\frac{144}{t}\right)\).
于是
\(S'(t)=\frac14\left(3t^2+24-\frac{144}{t^2}\right) =\frac{3(t^2-4)(t^2+12)}{4t^2}\).
所以 \(S(t)\) 在 \((0,2)\) 上递减，在 \((2,+\infty)\) 上递增，故当 \(t=2\) 时取最小值：
\(S_{\min}=S(2)=\frac{(4+12)^2}{8}=32\).
\end{solution}

\begin{problem}
已知椭圆 \(C:\dfrac{x^2}{a^2}+\dfrac{y^2}{b^2}=1\) 过点 \(A(-2,-1)\)，且 \(a=2b\).

\begin{enumerate}
    \item 求椭圆 \(C\) 的方程；
    \item 过点 \(B(-4,0)\) 的直线 \(l\) 交椭圆 \(C\) 于点 \(M\)，\(N\)，直线 \(MA\)，\(NA\) 分别交直线 \(x=-4\) 于点 \(P\)，\(Q\). 求 \(\left|\dfrac{PB}{BQ}\right|\) 的值.
\end{enumerate}

\end{problem}

\begin{solution}
（1）由题意得
\(\frac{(-2)^2}{a^2}+\frac{(-1)^2}{b^2}=1\)，\(a=2b\).
代入可得
\(\frac{4}{a^2}+\frac{1}{b^2}=1\)，
解得
\(a^2=8\)，\(b^2=2\).
所以椭圆方程为
\(\frac{x^2}{8}+\frac{y^2}{2}=1\).

（2）设 \(M(x_1,y_1)\)，\(N(x_2,y_2)\)，直线 \(MN\) 的方程为
\(y=k(x+4)\).
与椭圆方程联立可得
\((4k^2+1)x^2+32k^2x+64k^2-8=0\)，
于是
\(x_1+x_2=-\frac{32k^2}{4k^2+1}\)，\(x_1x_2=\frac{64k^2-8}{4k^2+1}\).
由直线 \(MA\)，\(NA\) 的方程以及 \(x=-4\) 可得点 \(P\)，\(Q\) 的纵坐标分别为
\(y_P=-\frac{(2k+1)(x_1+4)}{x_1+2}\)，\(y_Q=-\frac{(2k+1)(x_2+4)}{x_2+2}\).
进一步可算得
\(y_P+y_Q=0\)，
所以 \(y_P=-y_Q\). 因此
\(\left|\frac{PB}{BQ}\right|=\left|\frac{y_P}{y_Q}\right|=1\).
\end{solution}

\begin{problem}
已知 \(\{a_n\}\) 是无穷数列. 给出两个性质：

\begin{circlelist}
\item 对于 \(\{a_n\}\) 中任意两项 \(a_i\)，\(a_j\)（\(i>j\)），在 \(\{a_n\}\) 中都存在一项 \(a_m\)，使得 \(\dfrac{a_i^2}{a_j}=a_m\)；
\item 对于 \(\{a_n\}\) 中任意项 \(a_n\)（\(n\ge3\)），在 \(\{a_n\}\) 中都存在两项 \(a_k\)，\(a_l\)（\(k>l\)），使得 \(a_n=\dfrac{a_k^2}{a_l}\).
\end{circlelist}

\begin{enumerate}
\item 若 \(a_n=n\)（\(n=1\)，\(2\)，\(\dots\)），判断数列 \(\{a_n\}\) 是否满足性质（1），说明理由；
\item 若 \(a_n=2^{n-1}\)（\(n=1\)，\(2\)，\(\dots\)），判断数列 \(\{a_n\}\) 是否同时满足性质（1）和性质（2），说明理由；
\item 若 \(\{a_n\}\) 是递增数列，且同时满足性质（1）和性质（2），证明：\(\{a_n\}\) 为等比数列.
\end{enumerate}
\end{problem}

\begin{solution}
（1）当 \(a_n=n\) 时，取 \(a_2=2\)，\(a_3=3\)，则
\(\frac{a_3^2}{a_2}=\frac92\)，
故数列 \(\{a_n\}\) 不满足性质（1）.

（2）当 \(a_n=2^{n-1}\) 时，对任意 \(i>j\)，有
\(\frac{a_i^2}{a_j}=2^{2i-j-1}\)，
仍是数列中的某一项，所以满足性质（1）；对任意 \(n\ge3\)，取 \(k=n-1\)，\(l=n-2\)，则
\(\frac{a_k^2}{a_l}=2^{2(k-1)-(l-1)}=2^{n-1}=a_n\)，
所以也满足性质（2）.

（3）首先，性质（1）中的分式有意义，故 \(a_n\ne0\)（\(n\in\N^*\)）. 下面证明数列中的项同号.
若数列中同时含有正项和负项，设 \(N_0=\max\{n\mid a_n<0\}\).

\(1^{\odot}\) 若 \(N_0=1\)，即 \(a_1<0<a_2<a_3<\cdots\).
由性质（1）可知，存在 \(m_1\)，\(m_2\in\N^*\)，使得
\(a_{m_1}=\frac{a_2^2}{a_1}<0\)，\(a_{m_2}=\frac{a_3^2}{a_1}<0\).
由 \(N_0=1\) 可知，数列中唯一的负项是 \(a_1\)，从而 \(m_1=m_2=1\)，所以
\(\frac{a_2^2}{a_1}=\frac{a_3^2}{a_1}\)，
即 \(a_2=a_3\)，与数列的单调性矛盾，假设不成立.

\(2^{\odot}\) 若 \(N_0\ge2\)，由性质（1）知，存在 \(m\in\N^*\)，使得
\(a_m=\frac{a_{N_0}^2}{a_1}<0\).
由 \(N_0\) 的定义可知 \(m\le N_0\).
另一方面，由 \(a_1<a_{N_0}<0\) 得
\(a_m=\frac{a_{N_0}^2}{a_1}>\frac{a_{N_0}^2}{a_{N_0}}=a_{N_0}\)，
由数列的单调性可知 \(m>N_0\)，这与 \(m\le N_0\) 矛盾.

故数列中不可能同时含有正项和负项，从而数列中的项同号.

若对任意 \(n\in\N^*\) 都有 \(a_n<0\)，令 \(b_n=-a_n\)，则 \(\{b_n\}\) 是递减正数列.
对于任意 \(i>j\)，由性质（1）可知，在 \(\{a_n\}\) 中存在一项 \(a_m\)，使得
\(\frac{a_i^2}{a_j}=a_m\).
于是
\(\frac{b_i^2}{b_j}=\frac{(-a_i)^2}{-a_j}=-\frac{a_i^2}{a_j}=-a_m=b_m\)，
所以数列 \(\{b_n\}\) 也满足性质（1）.

对于任意 \(n\ge3\)，由性质（2）可知，在 \(\{a_n\}\) 中存在两项 \(a_k\)，\(a_l\)（\(k>l\)），使得
\(a_n=\frac{a_k^2}{a_l}\).
于是
\(b_n=-a_n=-\frac{a_k^2}{a_l}=\frac{(-a_k)^2}{-a_l}=\frac{b_k^2}{b_l}\)，
所以数列 \(\{b_n\}\) 也满足性质（2）.

下面分别处理全正和全负两种情形.

先设 \(a_n>0\). 由性质（2）取 \(n=3\)，存在 \(k>l\) 使得
\(a_3=\frac{a_k^2}{a_l}\).
因数列递增，\(a_k>a_l>0\)，故
\(a_3=a_k\cdot\frac{a_k}{a_l}>a_k\)，从而 \(k<3\)，只能有
\(k=2\)，\(l=1\). 因此 \(a_3=\frac{a_2^2}{a_1}\)，前三项成等比数列.

用数学归纳法证明全正情形下数列为等比数列. 假设前 \(k\)（\(k\ge3\)）项成等比数列，写成
\(a_s=a_1q^{s-1}\)（\(1\le s\le k\)），其中 \(q>1\).
由性质（1）取 \(i=k\)，\(j=k-1\)，存在 \(m\) 使得
\(a_m=\frac{a_k^2}{a_{k-1}}=a_1q^k>a_k\)，故 \(m\ge k+1\)，从而
\(a_1q^k=a_m\ge a_{k+1}\).
由性质（2）存在 \(s>t\) 使得 \(a_{k+1}=\frac{a_s^2}{a_t}\). 由于 \(a_s>a_t>0\)，\(a_{k+1}=a_s\frac{a_s}{a_t}>a_s\)，所以 \(a_s<a_{k+1}\)，从而 \(s\le k\)，\(t<s\le k\). 由归纳假设，
\(a_{k+1}=\frac{a_s^2}{a_t}=a_1q^{2s-t-1}\).
又 \(a_{k+1}>a_k=a_1q^{k-1}\)，所以
\(2s-t-1>k-1\). 结合 \(a_{k+1}\le a_1q^k\) 得
\(2s-t-1\le k\). 指数为整数，故 \(2s-t-1=k\)，即
\(a_{k+1}=a_1q^k\). 归纳完成.

再设所有 \(a_n<0\)，令 \(b_n=-a_n\). 上面已验证 \(\{b_n\}\) 仍满足性质（1）、（2），且是递减正数列.
由性质（2）取 \(n=3\)，存在 \(k>l\) 使得
\(b_3=\frac{b_k^2}{b_l}\). 因 \(b_k<b_l\)，有
\(b_3<b_k\)，递减性推出 \(k<3\)，故 \(k=2\)，\(l=1\)，即
\(b_3=\frac{b_2^2}{b_1}\)，前三项成等比数列.

用数学归纳法处理递减情形. 假设前 \(k\)（\(k\ge3\)）项成等比数列，写成
\(b_s=b_1q^{s-1}\)（\(1\le s\le k\)），其中 \(0<q<1\).
由性质（1）取 \(i=k\)，\(j=k-1\)，存在 \(m\) 使得
\(b_m=\frac{b_k^2}{b_{k-1}}=b_1q^k<b_k\)，故 \(m\ge k+1\)，从而
\(b_1q^k=b_m\le b_{k+1}\).
由性质（2）存在 \(s>t\) 使得 \(b_{k+1}=\frac{b_s^2}{b_t}\). 由于 \(0<b_s<b_t\)，\(b_{k+1}=b_s\frac{b_s}{b_t}<b_s\)，所以 \(b_s>b_{k+1}\)，递减性给出 \(s\le k\)，\(t<s\le k\). 由归纳假设，
\(b_{k+1}=\frac{b_s^2}{b_t}=b_1q^{2s-t-1}\).
又 \(b_{k+1}<b_k=b_1q^{k-1}\)，所以 \(2s-t-1>k-1\). 结合
\(b_{k+1}\ge b_1q^k\) 和 \(0<q<1\)，得 \(2s-t-1\le k\). 指数为整数，故
\(2s-t-1=k\)，即 \(b_{k+1}=b_1q^k\). 因而 \(\{b_n\}\)，从而 \(\{a_n\}\)，均为等比数列.
\end{solution}
